\subsection{A mesma assinatura vista de três lugares}
\label{sec:anel-tres-lugares}
\index{assinatura em anel@assinatura em anel!distribuicao das respostas@distribuição das respostas}

Há um pormenor curioso na conta precedente. Alice começou pelo segundo
lugar, e o verificador começou pelo primeiro. Em que lugar começaria
uma pessoa que recebesse a assinatura e procurasse descobrir como foi
feita? Vejamos as três possibilidades sobre a curva da
secção~\ref{sec:curva-137}. Usaremos o seu subgrupo de ordem 19, para
podermos publicar agora, juntamente com cada ponto, o escalar que
permite conferir a conta.

As três chaves públicas, conservadas nesta ordem, são
\[
\begin{array}{c|c|c}
 i&K_i&\text{coordenadas}\\ \hline
 1&3G &(23,41)\\
 2&5G &(56,60)\\
 3&11G&(93,33)
\end{array}
\]
Nesta assinatura didáctica conservemos a mensagem e o anel fixos.
Escrevamos \(h(L)\) para o desafio que o hash produz quando a ronda
termina no ponto \(L\). Eis os três valores de que vamos precisar.
\[
\begin{array}{c|c|c}
 L&\text{coordenadas}&h(L)\\ \hline
 8G &(44,33)&5\\
 4G &(96,53)&16\\
 14G&(81,60)&18
\end{array}
\]
Os valores correspondem a uma codificação inteiramente determinada.
Escreve-se cada ponto em dois bytes, abcissa seguida de ordenada.
Aos quatro bytes ASCII de \texttt{anel} seguem-se o byte de mensagem
\texttt{07}, os três pontos \(K_i\) e, por fim, \(L\). Interpreta-se
o resultado de Keccak--256 como inteiro \emph{little endian} e toma-se
o resto módulo 19. O quadro poupa ao leitor a avaliação do hash. As
coordenadas pequenas e o domínio escolhido pertencem a este exemplo.

Alice conhece a segunda chave privada, 5. Toma \(\alpha=8\) e escolhe
as duas respostas simuladas \(r_1=4\), \(r_3=6\). O primeiro valor
a colher é \(c_3=h(8G)=5\). Na terceira posição chega então ao ponto
\[
       L_3=6G+5(11G)=61G=4G,
       \qquad c_1=h(4G)=16 .
\]
O índice voltou ao princípio. Para a primeira chave, o ponto fica
\[
       L_1=4G+16(3G)=52G=14G,
       \qquad c_2=h(14G)=18 .
\]
Chegou-lhe o desafio de que precisava no lugar onde conhece o segredo.
A resposta é \(r_2=8-18\cdot5=-82\equiv13\pmod{19}\).
A assinatura completa que entrega é, por conseguinte,
\[
                (c_1,r_1,r_2,r_3)=(16,4,13,6).
\]
A disposição dos quatro números já não guarda a ordem temporal em que
Alice os descobriu. O primeiro a ser calculado, \(c_3\), nem sequer
figura nela.

Tomemos a folha tal como chega ao verificador. A primeira linha começa
por 16, e o que sai de cada linha fornece o desafio da próxima.
Podemos assim seguir o percurso inteiro na mesma folha.
\[
\begin{array}{c|r|r|l|r}
 i&c_i&r_i&L_i=r_iG+c_iK_i& h(L_i)\\ \hline
 1&16&4 &(4+16\cdot3)G=14G&18\\
 2&18&13&(13+18\cdot5)G=8G&5\\
 3&5&6  &(6+5\cdot11)G=4G&16
\end{array}
\]
Na segunda linha, \(13+90=103=5\cdot19+8\). O verificador
reencontrou assim o ponto de que Alice partira. Reconhece-se o fecho
pelo 16 da última casa, igual ao 16 da primeira.

Suponhamos agora que o dono da \emph{primeira} chave, 3, recebera a
mesma mensagem para assinar. Poderia ter produzido precisamente esta
folha? Se tivesse escolhido o escalar efémero 14 e as respostas
simuladas \(r_2=13,r_3=6\), partiria de \(L_1=14G\), obteria
\(c_2=18\), depois \(c_3=5\), depois \(c_1=16\). Ao regressar
ao seu lugar escreveria
\[
                     r_1=14-16\cdot3=-34\equiv4 .
\]
É a mesma assinatura, até ao último escalar. Falta o dono da chave
11. Com o escalar efémero 4 e as respostas simuladas
\(r_1=4,r_2=13\), partiria de \(L_3=4G\), obteria o desafio 16
para a primeira linha e acabaria por receber 5. Escreveria
\(r_3=4-5\cdot11=-51\equiv6\). A folha que entregaria seria,
também ela, a mesma.

Há aqui mais do que três maneiras possíveis de contar uma história.
Podemos comparar as suas probabilidades. Fixemos a mensagem, o anel,
a função \(h\) e esta particular assinatura. Para o dono da segunda
chave produzi-la, teve de escolher exactamente \((8,4,6)\) para
\((\alpha,r_1,r_3)\). As respostas simuladas lêem-se na assinatura,
e o desafio da segunda posição recompõe-se. Logo, o seu único
\(\alpha\) possível é \(13+18\cdot5\equiv8\).
Há exactamente um triplo de escolhas para o dono da primeira chave,
e um para o da terceira, pelas contas que acabámos de fazer.
Com escolhas uniformes independentes em \(\mathbb F_{19}\), cada
um desses triplos ocorre com a mesma probabilidade
\[
                          \frac{1}{19^3}.
\]
A probabilidade é a de \emph{produzir a folha fixada}, sob cada uma
das três hipóteses sobre o signatário. Se antes de a ver as três
posições tinham igual probabilidade, continuam a ter \(1/3\)
cada uma depois de a ver. Se o observador tinha razões exteriores
para atribuir probabilidades diferentes às três posições, multiplicar
as três pelo mesmo factor conserva as proporções iniciais.

A mesma conta vale para qualquer assinatura SAG aceite na construção
aqui definida. Para a posição \(i\), o escalar efémero compatível é
\(r_i+c_i k_i\), e as outras \(n-1\) respostas são precisamente as
que aparecem na assinatura. A probabilidade das escolhas é
\(l^{-n}\) para cada dono, fixados os restantes dados e admitindo
a amostragem uniforme em todo o corpo. Este argumento explica a
simetria do ciclo simples. Quando acrescentarmos uma imagem de chave,
haverá outro ponto público na declaração. A sua distribuição e a
dificuldade de o associar a uma chave pedem uma hipótese própria.
